\begin{statementbox}
	\textbf{\textcolor{CStatement}{条件}}
	\begin{description}[style=nextline, leftmargin=2.8em, labelsep=0.5em, itemsep=0.35em]
		\item[\textbf{\textcolor{CStatement}{条件1（分母非零约束）：}}] 对于任意角 $\alpha$，分式 $\frac{1-\cos\alpha}{2}$、$\frac{1+\cos\alpha}{2}$ 分母恒为非零；正切半角公式的分母 $1+\cos\alpha$ 和 $\sin\alpha$ 需保证非零。
		\item[\textbf{\textcolor{CStatement}{条件2（正切附加限制）：}}] 使用 $\tan\frac{\alpha}{2} = \frac{\sin\alpha}{1+\cos\alpha}$ 时 $\cos\alpha\neq -1$；使用 $\tan\frac{\alpha}{2} = \frac{1-\cos\alpha}{\sin\alpha}$ 时 $\sin\alpha\neq 0$。
	\end{description}

	\textbf{\textcolor{CStatement}{结论}}
	\begin{description}[style=nextline, leftmargin=2.8em, labelsep=0.5em, itemsep=0.45em]
		\item[\textbf{\textcolor{CStatement}{结论一（半角平方公式）：}}]
			\textbf{\textcolor{CStatement}{直观描述：}}
			利用一倍角的余弦表示半角的正弦平方与余弦平方。
			\[
				\sin^2\frac{\alpha}{2} = \frac{1-\cos\alpha}{2}.
			\]
			\[
				\cos^2\frac{\alpha}{2} = \frac{1+\cos\alpha}{2}.
			\]

		\item[\textbf{\textcolor{CStatement}{推广结论（正切半角形式）：}}]
			\textbf{\textcolor{CStatement}{直观描述：}}
			半角的正切可通过正弦与余弦的两种组合等价表达。
			\[
				\begin{aligned}
					\tan\frac{\alpha}{2} &= \frac{\sin\alpha}{1+\cos\alpha} \\
					&= \frac{1-\cos\alpha}{\sin\alpha}.
			\end{aligned}
			\]
	\end{description}

\end{statementbox}
